\documentclass{article}
\usepackage{amsmath}
\usepackage{amsthm}

\newtheorem{theorem}{Problem}

\begin{document}

\begin{theorem}
Three pencils and a jumbo eraser cost $\$1.24$. Five pencils and a jumbo eraser cost $\$1.82$. No prices include tax. In cents, what is the cost of a pencil?
\end{theorem}

\begin{proof}
Step 0: Step 1. I need to solve this problem to find the cost of a pencil in cents. The problem gives me two scenarios involving pencils and jumbo erasers, along with their total costs. Let's first understand what's being asked.
Step 1: Step 2. The problem states:

        - Three pencils and one jumbo eraser cost $1.24.

        - Five pencils and one jumbo eraser cost $1.82.

        And it asks for the cost of a pencil in cents.
Step 2: Step 3. Since the prices are given in dollars, I should convert them to cents to make calculations easier because the answer is required in cents.

        - $1.24 is equal to 124 cents.

        - $1.82 is equal to 182 cents.

        So, I can rephrase the problem as:

        - 3 pencils + 1 jumbo eraser = 124 cents

        - 5 pencils + 1 jumbo eraser = 182 cents
Step 3: Step 4. Now, I need to find the cost of one pencil. Let's denote:

        - Let p be the cost of one pencil in cents.

        - Let e be the cost of one jumbo eraser in cents.

        So, I can write the equations based on the given information:

        Equation 1: 3p + e = 124

        Equation 2: 5p + e = 182
Step 4: Step 5. I have a system of two equations with two variables (p and e). To find the value of p, I can use the method of elimination or substitution. I think elimination might be straightforward here.
Step 5: Step 6. If I subtract Equation 1 from Equation 2, the e (eraser cost) will cancel out, and I can solve for p.

        So, Equation 2 - Equation 1:

        (5p + e) - (3p + e) = 182 - 124

        Simplifying:

        5p + e - 3p - e = 58

        Which reduces to:

        2p = 58
Step 6: Step 7. Now, solving for p:

        p = 58 / 2

        p = 29 cents
Step 7: Step 8. So, the cost of one pencil is 29 cents. But to make sure, I should check if this makes sense with both original equations.
Step 8: Step 9. Plugging p = 29 into Equation 1:

        3(29) + e = 124

        87 + e = 124

        e = 124 - 87

        e = 37 cents
Step 9: Step 10. Now, plugging p = 29 and e = 37 into Equation 2 to verify:

        5(29) + 37 = 145 + 37 = 182 cents, which matches the second equation.
Step 10: Step 11. Everything checks out, so the cost of a pencil is indeed 29 cents.

        # Final Answer

        \[\boxed{29}\]
\end{proof}

\end{document}
